\section*{第 \ref{ch:b2:cell-signalling} 章——化学信使与信号转导}

\begin{solution}{exo:b2:cell-signalling:1}
胰岛素：内分泌，水溶性（肽）。皮质醇：内分泌，
脂溶性。突触处的乙酰胆碱：突触信号，水溶性。
伤口处的组胺：旁分泌，水溶性。雌二醇：
内分泌，脂溶性。来自肾上腺的肾上腺素：内分泌，
水溶性。一氧化氮：旁分泌，脂溶性（一种
可穿过膜的气体）。
\end{solution}

\begin{solution}{exo:b2:cell-signalling:2}
肾上腺素与受体结合（终止：解离，受体内化
）；受体激活\omterm{prop:b2:cell-signalling:gpcr}{G蛋白}，GDP换为GTP（终止：
GTP在数秒内水解）；\omterm{prop:b2:cell-signalling:gpcr}{G蛋白}激活腺苷酸环化酶，
生成cAMP（终止：磷酸二酯酶）；cAMP激活激酶A
（终止：cAMP消失）；激酶A使磷酸化酶激酶磷酸化，
后者使糖原磷酸化酶磷酸化（终止：磷酸酶）；
磷酸化酶从糖原中释放1-磷酸葡萄糖，转化
后以葡萄糖形式输出。
\end{solution}

\begin{solution}{exo:b2:cell-signalling:3}
表面受体在膜上结合水溶性信使，
以数秒的时间通过\omterm{prop:b2:cell-signalling:gpcr}{第二信使}和激酶起作用，改变
已有蛋白质的活性。\omterm{prop:b2:cell-signalling:nuclear}{核受体}在细胞内结合脂溶性
信使，以数小时的时间作为\omterm{prop:b2:cell-differentiation:transcriptional}{转录因子}起作用，
改变所合成的蛋白质种类。
\end{solution}

\begin{solution}{exo:b2:cell-signalling:4}
环 AMP：由腺苷酸环化酶合成，
由磷酸二酯酶分解。钙：经通道进入，或
从内质网释放，响应三磷酸肌醇，
由泵移除。三磷酸肌醇（与二酰甘油）：由
磷脂酶C作用于膜脂质生成，由磷酸酶移除。
\end{solution}

\begin{solution}{exo:b2:cell-signalling:5}
$\theta = [L]/(K_d + [L])$：$0.09$、$0.5$、$0.91$、$0.99$。受体数加倍
后占有率不变，而被占据受体数目加倍，
进而响应也加倍。
\end{solution}

\begin{solution}{exo:b2:cell-signalling:6}
$50\times 200\times 100\times 300 = 3\times 10^{8}$。100 个被占据的
受体每秒产生 $3\times 10^{10}$ 个产物分子——这是
一个上限，因为各阶段都趋于饱和。
\end{solution}

\begin{solution}{exo:b2:cell-signalling:7}
被锁定的\omterm{prop:b2:cell-signalling:gpcr}{G蛋白}使环化酶持续开启，cAMP保持高水平，激酶A
使肠道细胞的氯离子通道持续处于磷酸化状态，
即开放状态，氯离子被分泌，钠和水随之进入
肠腔——每天数升。这种锁定是对\omterm{prop:b2:cell-signalling:gpcr}{G蛋白}的共价修饰，
只有细胞本身被更新时才能解除，
数天之后才会发生。
\end{solution}

\begin{solution}{exo:b2:cell-signalling:8}
$0.9\times 10^{-6}\,\text{mol/L}\times 2\times 10^{-12}\,\text{L} =
1.8\times 10^{-18}$ 摩尔，$1.1\times 10^{6}$ 个离子（实际更多，因为
缓冲液结合了大部分进入的离子）。单个通道每秒运送 $10^{6}$ 个离子，
需一秒；一千个通道只需一毫秒，
这正是信号所需的时间尺度。
\end{solution}

\begin{solution}{exo:b2:cell-signalling:9}
两条通路汇聚于相同的酶：\omterm{prop:b2:cell-signalling:glucose}{胰高血糖素}的激酶A
使糖原合酶磷酸化（失活），并使磷酸化酶激酶
磷酸化（激活）；胰岛素的级联反应激活磷酸酶，以去除磷酸基团。
酶的状态取决于激酶与磷酸酶活性的平衡，
因此起决定作用的是两种\omterm{def:b2:cell-signalling:kinds}{激素}的比值；
两者均高时，通常磷酸酶占优，糖原
得以储存。
\end{solution}

\begin{solution}{exo:b2:cell-signalling:10}
肾上腺素：cAMP 在不到一秒内扩散穿过一个 \qty{20}{\micro m} 的细胞
（$x^{2}/2D$，其中 $D \approx \qty{3e-10}{m^2/s}$），而级联中的每一种
酶都已存在——只需数秒。皮质醇：一条
转录本需 \qty{60}{s}，蛋白质需 \qty{100}{s}；由
数十个 mRNA，每个携带十二个核糖体，细胞每分钟可制造
数百个酶分子，所以产生可测效应所需的数千个酶分子
需要数小时，发生在皮质醇进入细胞并抵达 DNA
所需数分钟之后。
\end{solution}

\begin{solution}{exo:b2:cell-signalling:11}
该受体持续发出信号：MAP 激酶级联驱动
细胞在没有任何生长因子的情况下分裂，并且忽略了
通常用于抑制它的信号缺失。由于单个
\omterm{def:b2:genome-diversification:mutation}{突变}即可去除对分裂的调控，此类受体（及其下游的
激酶）是最常见的癌基因之一。一种能占据激酶
ATP 结合位点的药物可阻止磷酸化，
并使受体失活——这正是多种抗癌药物的原理。
\end{solution}

\begin{solution}{exo:b2:cell-signalling:12}
内分泌：一分钟内到达每个细胞，逐个细胞成本低，无法
针对单个细胞，作用持续数分钟至数天——对
全身性状态（禁食、生长、生殖）不可或缺。神经系统：
毫秒级，靶点单一，布线成本高，作用持续
毫秒——对运动和快速反射不可或缺。两者都用于
速度与覆盖范围同时要求高的同一信息：交感神经与
肾上腺髓质都在惊吓时输送
去甲肾上腺素或肾上腺素。
\end{solution}

\begin{solution}{pb:b2:cell-signalling:1}
\textbf{1.} $1\,\text{nmol}/5\,\text{L} = \qty{0.2}{nmol/L}$。
\textbf{2.} $\theta = 0.2/1.2 = 0.17$：约 3300 个被占据的受体。
\textbf{3.} 三个半衰期：\qty{0.025}{nmol/L}；$\theta = 0.024$。
\textbf{4.} 约半分钟至一分钟，即\omterm{def:b2:blood-circulation:blood}{血液}
循环一周所需的时间；神经则在一秒之内把其
\omterm{def:b2:cell-signalling:kinds}{神经递质}直接送至肝脏。
\textbf{5.} 占有率相同，但被占据的受体为 $33\,000$ 个，
响应大十倍。
\textbf{6.} 在 \qty{1}{nmol/L} 时，一个 $K_d = \qty{1}{\micro
mol/L}$ 的受体被占据的比例仅为千分之一：无响应。令 $K_d$ 与
循环中的浓度范围相匹配，使占有率曲线陡峭的部分
恰好落在\omterm{def:b2:cell-signalling:kinds}{激素}实际变化的范围内。
\textbf{7.} $20\times 100\times 1\times 50\times 50\times 1000 =
5\times 10^{9}$ 个葡萄糖分子，每秒每个被占据的受体。
\textbf{8.} $3300\times 20\times 100 = 6.6\times 10^{6}$ 个 cAMP，
在第一秒内；在 \qty{5e-12}{L} 中，$1.1\times 10^{-17}$ 摩尔，约
\qty{2}{\micro mol/L}。
\textbf{9.} $3300\times 5\times 10^{9} = 1.7\times 10^{13}$ 每秒，$10^{15}$
每个细胞每分钟。
\textbf{10.} $2\times 10^{11}$ 个细胞 $\times 10^{15} = 2\times
10^{26}$ 个分子每分钟，330 摩尔，\qty{6e4}{g} 每分钟——为
观测值 \qty{5}{g/min} 的一万倍。增益相乘的前提是每一步都未饱和；
事实上磷酸化酶受其底物与磷酸酶的限制，
葡萄糖的输出又受转运蛋白的限制。
级联反应的增益用于换取速度与灵敏度，
而不是输出量。
\textbf{11.} 按 \qty{5}{g/min} 计，需二十分钟；惊吓在
数分钟内即结束，关闭机制早在此之前就已使级联反应停止。
\textbf{12.} G 蛋白对 GTP 的水解（数秒）；
磷酸二酯酶作用于 cAMP（数秒）；磷酸酶作用于
磷酸化的酶（数秒至数分钟）；\omterm{prop:b2:reproductive-hormones:pulses}{受体
脱敏}（数分钟）。最快的是 G 蛋白自身的“计时”
与磷酸二酯酶。
\textbf{13.} cAMP 并未被降解：它会累积并持续存在，
激酶持续激活，葡萄糖输出更大且持续更久
——这正是咖啡提神的作用。
\textbf{14.} $4000/50 = \qty{80}{s}$。
\textbf{15.} $600/5 = \qty{120}{s}$ 每个酶、每个核糖体；每个
核糖体每小时 30 个，每个 mRNA 每小时 600 个。
\textbf{16.} $100\times 600 = 60\,000$ 个酶每小时；$10^{5}$ 个酶
约需 \qty{1.7}{h}。
\textbf{17.} 皮质醇进入细胞、结合并到达其基因需十至二十分钟，
之后一小时，mRNA 才大量产生并被翻译，
再经两小时的积累：共约四小时。
\textbf{18.} 级联反应改变已存在的酶的活性；
而皮质醇的作用是改变哪些酶存在，持续数天，这只能
通过转录才能完成。惊吓不能等待四小时以获得新的
酶。
\textbf{19.} 受体增多意味着在任何肾上腺素浓度下，被占据的受体都更多：
占有率不变，但各浓度下的
响应更大——曲线
是整体放大而非平移。慢速\omterm{def:b2:cell-signalling:kinds}{激素}设定了快速\omterm{def:b2:cell-signalling:kinds}{激素}的增益。
\textbf{20.} 葡萄糖进入 $\beta$ 细胞并被代谢；ATP
升高，关闭 ATP 敏感的钾通道；膜
去极化；电压门控钙通道开放；钙离子
触发胰岛素颗粒的胞吐——代谢通过一个电学步骤
与分泌偶联。
\textbf{21.} 两种受体启动两条通路，它们汇聚于
相同的酶：\omterm{prop:b2:cell-signalling:glucose}{胰高血糖素}的激酶A使其磷酸化，胰岛素的
级联反应激活使其去磷酸化的磷酸酶；酶的
状态，进而糖原代谢的方向，取决于
两者的比值。
\textbf{22.} 空腹血糖偏高（肝脏不受抑制，持续
生成葡萄糖），餐后血糖极高（肌肉与脂肪
不摄取葡萄糖）；无胰岛素时，脂肪储备释放脂肪酸，肝脏将其
转化为酮酸——即酮症酸中毒。
\textbf{23.} 运动中的肌肉无需胰岛素即可摄取葡萄糖，因此在
长时间运动中血糖下降，而注射的胰岛素不受调节，
仍持续抑制肝脏：低血糖与虚脱。
\textbf{24.} 血糖回路：感受器为 $\beta$ 与 $\alpha$ 细胞，
两种效应\omterm{def:b2:cell-signalling:kinds}{激素}分别作用于两个方向，高增益使
水平保持在五分之一以内，作用需数分钟。\omterm{def:b2:blood-pressure:baroreflex}{压力反射}：牵张
感受器，以\omterm{def:b2:heart:anatomy}{心脏}和血管为效应器，增益约为四，作用
需一秒。
\textbf{25.} 释放后占有率为 $0.17$；增益为每秒 $5\times 10^{9}$
（每个被占据的受体）；肝脏实际输出 \qty{5}{g/min}；皮质醇
需约四小时。
\end{solution}
